\documentclass[acmsmall,review,screen]{acmart}

\usepackage{booktabs}
\usepackage{tabularx}
\usepackage{array}
\usepackage{multirow}
\usepackage{xspace}
\usepackage{tikz}
\usepackage{pgfplots}
\usetikzlibrary{arrows.meta,positioning,calc,fit}
\pgfplotsset{compat=1.18}
\setlength{\headheight}{19pt}

\newcommand{\method}{ARCIS\xspace}
\newcolumntype{Y}{>{\raggedright\arraybackslash}X}

\setcopyright{none}
\settopmatter{printfolios=true}
\acmJournal{TOMM}
\acmYear{2026}
\acmMonth{9}

\begin{document}

\title{ARCIS: Authenticated Robust Cover-Selection Image Signaling under Finite-Index Constraints}

\author{Arthur Ulrich Ewane}
\affiliation{\institution{Department of Computer Science, University of Yaoundé I, LIRIMA, TEAM GRIMCAPE}\city{Yaoundé}\country{Cameroon}}
\email{ulrich.ewane@facsciences-uy1.cm}

\author{Juvet Karnel Sadié}
\affiliation{\institution{University of Yaoundé I / UMMISCO, IRD France Nord / Sorbonne Université, UMI 209 UMMISCO}\city{Yaoundé / Bondy / Paris}\country{Cameroon / France}}
\email{sadie.juvet@gmail.com}

\author{Stéphane Gaël R. Ekodeck}
\affiliation{\institution{University of Yaoundé I / UMMISCO, IRD France Nord / Sorbonne Université, UMI 209 UMMISCO}\city{Yaoundé / Bondy / Paris}\country{Cameroon / France}}
\email{stephane-gael.ekodeck@facsciences-uy1.cm}

\author{Serge Alain Ebélé}
\affiliation{\institution{University of Yaoundé I / UMMISCO, IRD France Nord / Sorbonne Université, UMI 209 UMMISCO}\city{Yaoundé / Bondy / Paris}\country{Cameroon / France}}
\email{serge-alain.ebele@facsciences-uy1.cm}

\author{Vivien Loïck Beyala Kamgang}
\affiliation{\institution{The University of Bertoua / RUFCEA (UR-IFIA), University of Dschang}\city{Bertoua / Dschang}\country{Cameroon}}
\email{vivsxx@univ-bertoua.cm}

\author{Damaris Belle M. Fotso}
\affiliation{\institution{Department of Computer Science, University of Yaoundé I, LIRIMA}\city{Yaoundé}\country{Cameroon}}
\email{makwate.fotso@facsciences-uy1.cm}

\author{Chantal Marguerite Mveh-Abia}
\authornote{Corresponding author.}
\affiliation{\institution{University of Yaoundé I: Department of Computer Science, LIRIMA, TEAM GRIMCAPE; Department of Computer Engineering, National Advanced School of Engineering}\city{Yaoundé}\country{Cameroon}}
\email{cmmveh@yahoo.fr}

\begin{abstract}
Selection-based coverless communication preserves carrier pixels, but reliable transmission from a finite shared index depends jointly on symbol demand, channel robustness, coding overhead, and the distribution induced by cover selection. We present \method (Authenticated Robust Cover-Selection Image Signaling), a protocol that combines a learned normalized descriptor, calibration-locked quantile codebooks, complementary five-cover group banks, exact session-balanced keyed Gray mapping, AES-GCM authenticated framing, replay control, and Reed--Solomon (RS) correction. The final operating point ($K=8$, group size 5, RS128) was fixed on BOSSBase development evidence and then frozen for five Caltech-101 partitions. External calibration recovered 1,800/1,800 authenticated messages; eight unseen holdouts recovered 1,163/1,200 (96.92\%), with all 37 failures confined to the 12\% crop. A controlled RS-parity ablation, holding the representation, codebook, and group banks fixed, increased holdout recovery from 829/1,200 (69.08\%) without RS parity to 1,163/1,200 with RS128, while mean transmitted images per session increased from 970 to 2,676.7. At RS128, 8-, 32-, and 64-byte payloads require 2,320, 2,640, and 3,070 transmitted images, respectively. Under an explicitly conditional i.i.d.-uniform symbol-demand planning model, exact finite-index blocking probabilities for those three payload sizes are $5.04\times10^{-43}$, $1.64\times10^{-34}$, and $1.14\times10^{-25}$. A repeated image-disjoint detector audit yields mean macro-AUC 0.504 for SRM-lite/ExtraTrees, 0.513 for GLCM/logistic regression, and 0.527 for a 30-head CNN. The resulting evidence identifies both the reliability gains and the communication cost of robust finite-index signaling while bounding the security claim to the evaluated channel and warden models.
\end{abstract}

\begin{CCSXML}
<ccs2012>
 <concept><concept_id>10002978.10003006</concept_id><concept_desc>Security and privacy~Cryptography</concept_desc><concept_significance>500</concept_significance></concept>
 <concept><concept_id>10002951.10003227</concept_id><concept_desc>Information systems~Multimedia information systems</concept_desc><concept_significance>300</concept_significance></concept>
 <concept><concept_id>10010147.10010257.10010293</concept_id><concept_desc>Computing methodologies~Neural networks</concept_desc><concept_significance>300</concept_significance></concept>
</ccs2012>
\end{CCSXML}
\ccsdesc[500]{Security and privacy~Cryptography}
\ccsdesc[300]{Information systems~Multimedia information systems}
\ccsdesc[300]{Computing methodologies~Neural networks}

\keywords{coverless image steganography, cover selection, covert communication, multimedia security, authenticated encryption, finite-index feasibility, error correction, steganalysis}

\maketitle

\section{Introduction}
Image steganography traditionally hides information by modifying a carrier, whereas coverless approaches communicate through selection, indexing, semantic mapping, or generation. Selection preserves the original pixels, shifting the principal security question toward the distribution of transmitted images and the reliability of their labels after channel perturbations. Recent surveys accordingly separate cover-selection design from detectability and emphasize evaluation criteria that reflect the selected-carrier distribution~\cite{azam2025review,kombrink2025survey}.

A finite shared index adds a second protocol constraint. Every protected message induces a concrete multiplicity demand over codebook symbols. Robust communication therefore requires enough qualified carrier groups for every requested symbol across the complete protected envelope. Stability-aware quantization and representative selection improve robust cover choice~\cite{guo2026capacity}, while restoration and deep-hashing methods strengthen geometric robustness~\cite{meng2025restoration,tan2026hashing}. \method complements these directions by placing finite-index feasibility, authenticated message recovery, and selection detectability inside one end-to-end evaluation unit.

\method couples a channel-aware quantile codebook with complementary five-cover groups whose majority-failure signatures are characterized on calibration transformations. A proportional-deficit scheduler preserves the natural signature mixture of each label bank, while an exact cyclic keyed Gray mapping balances coded symbols across visual clusters over every full $K$-session cycle. AES-GCM binds payload and session state, Reed--Solomon coding absorbs residual byte errors, and authenticated plaintext recovery becomes the final success criterion. The scientific contribution is this integrated finite-index construction and its externally frozen evaluation rather than any individual standard primitive.

The validation separates development evidence from external evidence. BOSSBase 1.01~\cite{bas2011boss} supports design and operating-point selection. The final protocol then evaluates Caltech-101~\cite{caltech101,feifei2004caltech} with $K=8$, five covers per coded symbol, 128 Reed--Solomon parity bytes, and 30 authenticated sessions per partition. The external campaign combines twelve calibration transformations, eight unseen holdout transformations, complete-index usage measurements, a controlled RS-parity ablation, a finite-index blocking-risk analysis, and repeated selection-steganalysis with SRM-lite, GLCM, and CNN detectors.

The resulting evidence spans robustness, cost, feasibility, usage, and detectability. The BOSSBase development lineage achieves 1,800/1,800 calibration recoveries over five seeds, with 7,995--8,000 unique covers exercised per seed and repeated-session detector AUCs close to chance. On external Caltech partitions, calibration reaches 1,800/1,800 authenticated recoveries and exercises all 7,000 covers in every partition. Unseen holdouts reach 1,163/1,200 authenticated recoveries (96.92\%), with the strongest 12\% crop defining the observed channel boundary. The parity sweep quantifies the reliability--communication tradeoff rather than treating RS128 as cost-free. Final mean macro-AUCs are 0.504 for SRM-lite, 0.513 for GLCM, and 0.527 for the CNN.

The article does not claim universal undetectability, universal robustness, or direct numerical superiority over generative steganography. Its claims are bounded to the declared finite-index protocol, channel families, external corpus, and evaluated wardens.

\section{Related Work and Positioning}
\subsection{Selection, retrieval, and robust coverless communication}
Early coverless systems mapped image content or object configurations to symbols without modifying carriers. Multi-object recognition illustrates semantic cover selection~\cite{luo2021multiobject}. Recent reviews organize cover-selection strategies and their security implications~\cite{azam2025review,kombrink2025survey}. Guo and Ping~\cite{guo2026capacity} use stability-aware pseudo-Zernike quantization and representative selection to balance robustness and nominal capacity. \method focuses on a complementary question: whether the complete authenticated coded demand can be realized by a fixed qualified shared index and recovered at the plaintext level.

Geometric perturbations remain a central challenge. Meng et al.~\cite{meng2025restoration} restore attacked images before the underlying coverless decoder, while Tan et al.~\cite{tan2026hashing} use deep unsupervised hashing to improve resistance to geometric attacks. \method incorporates calibration transformations directly into codebook and group-bank construction, while eight stronger or intermediate transformation values remain holdouts. The observed 12\% crop boundary is therefore reported as an explicit external-validity result.

\subsection{Generative steganography and multimedia security}
Generative approaches communicate through synthesized content. DiffStega reconstructs a secret image through prompt- and password-conditioned diffusion~\cite{yang2024diffstega}. Guidance-feature-distribution steganography embeds information in a generative feature distribution~\cite{sun2024guidance}, while progressive generative steganography in TOMM studies capacity, extraction, and detectability during high-resolution generation~\cite{zhou2025progressive}. \method instead operates on an indexed natural-image carrier set and measures exact authenticated byte recovery. Each family is therefore reported with its native information object and metrics.
